\begin{onehalfspace}

Il Capitolo~\ref{chap:sviluppo} ha descritto M1 (TOP-1), M2
(classificatore TOP-1 seguito da TOP-4) e l'ablazione
$M_{\mathrm{TOP4}}$. Questa prima parte del capitolo presenta la campagna rigenerata
dopo la correzione dell'operatore $\times7\bmod15$ e dei contratti
statistici. Il risultato principale è che la ricerca multi-candidato
riduce le iterazioni, mentre il classificatore a valle non aggiunge
beneficio rispetto alla stessa ricerca priva di filtro.

% -----------------------------------------------
\section{Setup Sperimentale}
\label{sec:setup_m1}
\label{sec:setup_m2}

Il confronto usa $N=15$, $a=7$, con due preset di rumore uniformi e illustrativi. In ogni replica, le tre strategie ricevono gli stessi istogrammi. Le differenze dipendono quindi dalla regola di post-processing e non da campioni casuali diversi. Il modello include depolarizzazione, rilassamento, perdita di fase ed errori di lettura; non riproduce la calibrazione o la topologia di una QPU specifica. Ogni confronto principale comprende 30 repliche da $1\,024$ shot per iterazione. Si usa il test Wilcoxon--Pratt con correzione di Holm per i confronti multipli.

Il contratto completo --- circuito e versioni, seed, manifest, gestione
dei pareggi e della censura, insieme al resoconto della campagna storica
non riutilizzata --- è conservato nella
Sezione~\ref{app:fase1_contratto} dell'Appendice~\ref{app:fase1}.

% -----------------------------------------------
\section{Perimetro Ideale e Rumoroso}
\label{sec:circuito_m1}
\label{sec:scalabilita}

Le istanze maggiori servono prima di tutto a verificare la correttezza ideale. Per $N=21$, $a=2$, si controllano la tabella dell'operatore, l'azzeramento delle ancille e la distribuzione QPE. La compilazione raggiunge però 21\,036 CX e profondità 23\,081: questa istanza resta fuori dalla campagna rumorosa principale. Anche $N=35$, $a=6$, $r=2$ supera i controlli ideali. La distanza di variazione totale dalla distribuzione QPE teorica è $1.38\times10^{-9}$ per $N=21$ e $2.50\times10^{-13}$ per $N=35$. In entrambi i casi le ancille tornano a zero. Il proxy di nessun evento e la tabella del perimetro sono nella Sezione~\ref{app:fase1_contratto}. I confronti M1, TOP-4 e M2 che seguono riguardano $N=15$.

% -----------------------------------------------
\section{La Baseline: Metodo 1 (TOP-1)}
\label{sec:risultati_uc_m1}
\label{sec:analisi_m1}
\label{sec:uc1_m1}
\label{sec:uc2_m1}

La diagnostica indipendente di 60 repliche trova la moda $y=0$ in
$30/60$ casi e una moda utile negli altri $30/60$, mentre TOP-4 trova
sempre almeno un candidato utile. L'instabilità di TOP-1 deriva quindi da
quale dei quattro picchi teorici risulta dominante, non dalla completa
scomparsa del segnale. Gli istogrammi estremi, selezionati con criterio
prefissato e seed registrati, sono mostrati nella
Figura~\ref{fig:istogrammi_qpe} del presente capitolo.

% -----------------------------------------------
\section{L'Ipotesi ML e il Verdetto dell'Ablazione}
\label{sec:confronto}
\label{sec:uc1_m2}
\label{sec:uc2_m2}

I classificatori TOP-1 sono stati rigenerati su $2\,000$ istogrammi per scenario. Con TOP-16, tutti i $2\,000/2\,000$ campioni di entrambi gli scenari sono positivi: il controllo produce quindi una sola classe. Fra Random Forest, SVM e MLP viene scelta la SVM. Sul test indipendente di 400 istogrammi per scenario ottiene F1 $0.919/0.923$, accuratezza $0.895/0.888$ e AUC $0.965/0.958$ (UC1/UC2). La SVM predice bene se TOP-1 avrà successo. Per stabilire se sia utile, occorre però confrontare il costo di M2 con quello della ricerca senza filtro.

L'audit combinatorio spiega inoltre perché TOP-16 sia poco discriminante:
63 dei 256 esiti sono riconosciuti come utili dalla verifica classica
adottata. Se si scelgono candidati senza informazione, la probabilità di
includerne almeno uno è già circa il $99.07\%$ con 16 candidati e il
$67.94\%$ con quattro. Il successo TOP-K va quindi interpretato insieme
alla massa sui picchi, non come misura della fedeltà quantistica.
Metriche complete, protocollo train/validation/test e derivazione del
controllo combinatorio sono nelle Sezioni~\ref{sec:valutazione_clf}
e~\ref{subsec:soglia_classe_negativa}. La
Tabella~\ref{tab:riepilogo_rho} riporta invece la misura operativa,
calcolata sui medesimi istogrammi per tutte le strategie.

\begin{table}[H]
\centering
\caption{Confronto appaiato M1, TOP-4 e M2 ($30$ repliche per scenario)}
\label{tab:riepilogo_m1}
\label{tab:riepilogo_rho}
\begin{tabular}{llcccc}
\toprule
\textbf{Scenario} & \textbf{Strategia} & \textbf{Successo} &
$\boldsymbol{\bar M}$ & $\boldsymbol{\rho}$ vs M1 &
\textbf{$p_{\mathrm{Holm}}$ vs M1} \\
\midrule
UC1 & M1 & $30/30$ & 1.37 & --- & --- \\
UC1 & TOP-4 & $30/30$ & 1.00 & 1.367 & 0.0024 \\
UC1 & M2 & $30/30$ & 1.50 & 0.911 & 0.7314 \\
\midrule
UC2 & M1 & $30/30$ & 1.50 & --- & --- \\
UC2 & TOP-4 & $30/30$ & 1.00 & 1.500 & 0.0015 \\
UC2 & M2 & $30/30$ & 1.37 & 1.098 & 0.0874 \\
\bottomrule
\end{tabular}
\end{table}

TOP-4 converge al primo tentativo in tutte le repliche. M2 non migliora
significativamente M1 e, rispetto alla propria ablazione, richiede più
iterazioni: Wilcoxon bilaterale con correzione di Holm
$p=0.0033$ (UC1) e $p=0.0096$ (UC2). Il filtro scarta dunque istogrammi
che la ricerca TOP-4 avrebbe risolto; il guadagno appartiene alla ricerca
multi-candidato, non al classificatore.

Il confronto con la letteratura chiarisce il ruolo dell'ablazione. Yang e Markidis addestrano alberi e foreste casuali su $680$ distribuzioni raccolte da sistemi \ac{IBM} reali. Il modello predice se le frazioni continue recupereranno l'ordine vero. Usa quattro caratteristiche dell'istogramma, fra cui la frazione di massa verificata dominante, che risulta la più informativa~\cite{yang2026}. Quei dati provengono da hardware, mentre qui si usa un modello di rumore. Resta però una domanda diversa: il classificatore migliora una regola semplice applicata agli stessi istogrammi? L'ablazione di questa tesi misura proprio questo contributo e non trova un beneficio. Il risultato non smentisce la capacità di classificare gli output. Mostra che buone predizioni non bastano a dimostrare un vantaggio rispetto alla ricerca di più candidati.

\begin{figure}[H]
    \centering
    \includegraphics[width=0.75\textwidth]{fig_iterazioni_m1_m2.pdf}
    \caption[Analisi di ablazione M1, TOP-4 e M2]{Distribuzione delle
    iterazioni nelle 30 repliche appaiate. I fallimenti, se presenti,
    sarebbero mostrati separatamente con sentinella 51; in questa
    baseline tutte le repliche hanno successo.}
    \label{fig:iterazioni_m1_m2}
\end{figure}

% -----------------------------------------------
\section{Robustezza Esplorativa della Strategia TOP-4}
\label{sec:parametri_mitigazione}
\label{subsec:sintesi_parametri}

Otto famiglie di sweep variano un parametro per volta. I test sono
appaiati all'interno del singolo punto, ma i $p$ non ricevono una
correzione globale fra le otto famiglie: i risultati vanno quindi letti
come esplorativi. La Tabella~\ref{tab:sintesi_parametri} evita di
trasformare l'assenza di variazione osservata in una pretesa universale
di insensibilità al rumore.

\begin{table}[H]
\centering
\caption{Sintesi degli otto sweep schema~2 su $N=15$}
\label{tab:sintesi_parametri}
\begin{tabularx}{\textwidth}{llX}
\toprule
\textbf{Famiglia} & \textbf{Intervallo} & \textbf{Esito TOP-4} \\
\midrule
$K$ & $1,2,3,4,6,8$ & soglia osservata $K=2$ \\
$\lambda_{2q}$ & $10^{-3}$--$10^{-1}$ & $\bar M=1.00$--$1.07$, successo $100\%$ \\
shot & $128$--$2\,048$ & $\bar M=1.00$ \\
$K\times\lambda_{2q}$ & $K=2$--8, $5\times10^{-3}$--$5\times10^{-2}$ & $\bar M=1.00$ \\
$T_1/T_2$ & $20/16$--$500/400\,\mu$s & $\bar M=1.00$ \\
$\lambda_{1q}$ & $10^{-4}$--$2\times10^{-2}$ & $\bar M=1.00$ \\
readout simmetrico & $0$--$20\%$ & $\bar M=1.00$ \\
ottimizzazione & livelli 0--3 & $\bar M=1.00$ \\
\bottomrule
\end{tabularx}
\end{table}

Variando $K$, si raggiunge già una sola iterazione con $K=2$. I valori $K=3,4,6,8$ non migliorano questo risultato. TOP-4 funziona bene ed è coerente con i quattro picchi dell'istanza, ma non è un ottimo universale. Nei test non è nemmeno il numero minimo di candidati sufficiente a raggiungere una sola iterazione (Tabella~\ref{tab:sweep_k}).

Il livello di ottimizzazione modifica il circuito: livello 0 produce 242
CX e profondità 427, livello 1 produce 236/421, mentre i livelli 2 e 3
coincidono a 224/412. Nonostante ciò TOP-4 resta a una iterazione nei
quattro punti. Le oscillazioni di M1 negli sweep non sono monotone e, con
30 repliche, non sono interpretate come leggi causali.

\paragraph{Diagnostiche e confronto con ZNE.}
\label{subsec:frazione_coerente}
\label{subsec:diagnostiche_qpe}
La sola iterazione di successo non descrive quanta struttura periodica resti
nell'istogramma. La Figura~\ref{fig:istogrammi_qpe} mostra, fra le 30 repliche
nominali UC1, i casi con massa massima e minima sui quattro picchi teorici
$\{0,64,128,192\}$.

\begin{figure}[!htbp]
    \centering
    \includegraphics[width=0.88\textwidth]{fig_istogrammi_qpe.pdf}
    \caption{Repliche UC1 con massa massima e minima sui quattro picchi QPE
    teorici.}
    \label{fig:istogrammi_qpe}
\end{figure}

La massa media $m$ sui picchi viene convertita nell'indice descrittivo
\begin{equation}
    f=\frac{m-4/256}{1-4/256},
    \label{eq:miscela_coerente}
\end{equation}
che vale zero per il fondo uniforme e uno per una distribuzione concentrata
sui picchi. Non è una fedeltà. La Tabella~\ref{tab:frazione_coerente} lo
confronta con il proxy $P_0=(1-15\lambda_{2q}/16)^{224}$.

\begin{table}[!htbp]
\centering
\caption{Massa efficace sui picchi e proxy di nessun evento 2Q.}
\label{tab:frazione_coerente}
\begin{tabular}{ccc}
\toprule
$\lambda_{2q}$ & $f$ efficace & $P_0$ proxy \\
\midrule
$10^{-3}$ & 0.8036 & $8.11\times10^{-1}$ \\
$5\times10^{-3}$ & 0.7039 & $3.49\times10^{-1}$ \\
$10^{-2}$ & 0.5938 & $1.21\times10^{-1}$ \\
$2\times10^{-2}$ & 0.4286 & $1.44\times10^{-2}$ \\
$5\times10^{-2}$ & 0.1688 & $2.14\times10^{-5}$ \\
$10^{-1}$ & 0.0421 & $2.65\times10^{-10}$ \\
$2\times10^{-1}$ & 0.0047 & $6.32\times10^{-21}$ \\
$3\times10^{-1}$ & 0.0015 & $7.47\times10^{-33}$ \\
\bottomrule
\end{tabular}
\end{table}

\begin{figure}[!htbp]
    \centering
    \includegraphics[width=0.76\textwidth]{frazione_coerente.pdf}
    \caption{Il proxy di nessun evento decade più rapidamente della massa
    strutturata osservata sui picchi.}
    \label{fig:frazione_coerente}
\end{figure}

La struttura QPE sopravvive dunque anche quando il proxy moltiplicativo di
nessun evento Pauli è già prossimo a zero: i due indicatori non sono né
fedeltà né probabilità di fattorizzazione e non vanno identificati. Il
protocollo completo e gli altri sweep restano nell'Appendice~\ref{app:fase1}.

Nello stesso perimetro, la ZNE digitale non migliora M1 e porta il costo
medio da $1\,399$ a $3\,072$--$4\,813$ shot, mentre TOP-4 raggiunge
$\bar M=1.00$ con $1\,024$ shot. Il protocollo, che scala solo i parametri
depolarizzanti e non equivale al gate folding su hardware, e la tabella
appaiata completa sono conservati nella
Sezione~\ref{sec:app_zne} e nella Tabella~\ref{tab:confronto_zne}.
Il verdetto è locale al modello uniforme e a questa istanza, non una
valutazione generale della ZNE.

% -----------------------------------------------
\section{Bilancio della Verifica e Ponte verso la Correzione d'Errore}
\label{sec:top4_mitigazione}
\label{sec:validita}
\label{sec:confronto_zne}

La campagna consente quattro conclusioni delimitate dai dati:
\begin{enumerate}
    \item l'operatore $\times7\bmod15$ e il circuito compilato sono
    fissati da truth table, manifest e hash;
    \item TOP-4 riduce $\bar M$ da $1.37/1.50$ a $1.00$ nei due scenari;
    \item il classificatore TOP-1 non migliora M1 e peggiora la propria
    ablazione TOP-4;
    \item la robustezza osservata negli sweep è esplorativa e non rende
    il modello uniforme equivalente a una QPU reale.
\end{enumerate}

Il post-processing recupera informazione ancora presente
nell'istogramma, ma non può correggere uno stato già corrotto durante il
calcolo. Questo limite motiva il passaggio ai codici di correzione e alla
decodifica della sindrome nelle sezioni successive.

\end{onehalfspace}
